% ===================================================================
%  ТАБЛИЦА № 15 — Рынок. — Съестные припасы.
%  Traduction 2026 d'apres texto/io, controlee sur texto/fr, texto/en
%  et texto/es. Consigne : voir texto/ru/10-tablica-01.tex.
% ===================================================================

%%K t15-apar-1 apar
\VUsaut{4.0mm}
\VUtitre{15.0pt}{60}{ТАБЛИЦА № 15}
\VUsaut{3.0mm}

%%K t15-tit-1 sub
\VUpk{11.4pt}{40}{Рынок. --- Съестные припасы.}
\VUsaut{3.0mm}

%%K t15-01-1 p
1. --- Бо́льшая часть торговцев, которые снабжают нас нужными припасами,
собрана под \VUgras{навесом} \textsuperscript{(1)} \VUgras{крытого рынка} или на
соседних улицах.

%%K t15-02-1 p
2. --- \VUgras{Колбасник} \textsuperscript{(2)}, который продаёт \VUgras{окорок}
\textsuperscript{(3)}, \VUgras{кровяную колбасу} \textsuperscript{(4)},
\VUgras{колбасу} \textsuperscript{(5)}, \VUgras{сардельки} \textsuperscript{(6)} и
\VUgras{сало} \textsuperscript{(7)}, стоит рядом с \VUgras{торговкой маслом и
сыром} \textsuperscript{(8)}, чей прилавок уставлен \VUgras{голландскими сырами}
\textsuperscript{(9)} с красной коркой, \VUgras{камамберами} \textsuperscript{(10)},
большими \VUgras{кругами грюйера} \textsuperscript{(11)} и свежими
\VUgras{кусками масла} \textsuperscript{(12)}.

%%K t15-03-1 p
3. --- Перед нею \VUgras{зеленщица} \textsuperscript{(13)} предлагает своим
покупателям прекрасные \VUgras{помидоры} \textsuperscript{(14)},
\VUgras{цветную капусту} \textsuperscript{(15)}, \VUgras{салат} \textsuperscript{(16)},
\VUgras{капусту} \textsuperscript{(17)}, \VUgras{тыквы} \textsuperscript{(19)},
\VUgras{дыни} \textsuperscript{(18)} и даже \VUgras{грибы} \textsuperscript{(20)}. Но
\VUgras{пивовар}, который остановил свою \VUgras{развозную телегу}
\textsuperscript{(21)}, гружённую \VUgras{пивными бочонками} \textsuperscript{(22)},
прямо перед её прилавком, мешает покупателям подойти. Огородник несёт ей
корзину \VUgras{артишоков} \textsuperscript{(23)}.

%%K t15-tit-2 sub
\VUsaut{4.0mm}
\VUpk{10.2pt}{40}{Продавать и покупать.}
\VUsaut{3.0mm}

%%K t15-04-1 p
4. --- На рынке большое оживление, потому что настал час \VUgras{торгов}
\textsuperscript{(24)}. \VUgras{Хозяйки} \textsuperscript{(26)}, которые приходят сюда
за покупками, останавливаются по дороге перед прилавком
\VUgras{торговки потрохами} \textsuperscript{(25)}, чтобы купить \VUgras{рубец}
или \VUgras{телячью голову}.

%%K t15-05-1 p
5. --- Толстая \VUgras{кухарка} \textsuperscript{(27)} торгуется о прекрасном
\VUgras{тунце} с \VUgras{рыбницей} \textsuperscript{(28)}, которая набавляет цену
как ей вздумается. Ей привезли много \VUgras{угрей, камбалы, форели} и
\VUgras{карпов}. Её \VUgras{треска} и её \VUgras{мидии} очень свежи.

%%K t15-tit-3 sub
\VUsaut{4.0mm}
\VUpk{10.2pt}{40}{Дешёвое и дорогое.}
\VUsaut{3.0mm}

%%K t15-06-1 p
6. --- \VUgras{Торговка птицей} \textsuperscript{(29)}, устроившаяся рядом с нею,
очень честна. Продавая \VUgras{индюка}, \VUgras{гуся}, \VUgras{утку} или
простую \VUgras{курицу}, она не просит больше справедливой цены.

%%K t15-07-1 p
7. --- На углу площади, на первом этаже, недавно завела торговлю
\VUgras{бельевщица} \textsuperscript{(30)}, у которой покупатели всегда наверняка
найдут очень хороший выбор \VUgras{скатертей}, \VUgras{салфеток, фартуков}
и \VUgras{кухонных полотенец}. На том же этаже устроил свою мастерскую
\VUgras{ножовщик} \textsuperscript{(31)}. Он точит \VUgras{ножи} торговок и делает
огородникам \VUgras{садовые ножи} из самой твёрдой \VUgras{стали}.

%%K t15-tit-4 sub
\VUsaut{4.0mm}
\VUpk{10.2pt}{40}{Бакалея.}
\VUsaut{3.0mm}

%%K t15-08-1 p
8. --- Под его мастерской не так давно открылась большая съестная лавка.
Это уже не простая, скромная \VUgras{бакалейная лавочка} \textsuperscript{(32)}
былых времён, где продавали только обиходные товары --- \VUgras{чай}
\textsuperscript{(33)}, \VUgras{шоколад} \textsuperscript{(34)}, \VUgras{сахарную голову}
\textsuperscript{(37)} и \VUgras{мыло} \textsuperscript{(38)}. Здесь продают всё:
\VUgras{хрустящее печенье}, \VUgras{консервы} \textsuperscript{(35)},
\VUgras{ликёры} \textsuperscript{(36)}, \VUgras{ром, коньяк, кирш, анисовку} и
\VUgras{кюрасао}. Дичь достать здесь так же легко --- \VUgras{зайца}
\textsuperscript{(39)}, \VUgras{дроздов} \textsuperscript{(40)}, \VUgras{куропаток}
\textsuperscript{(41)}, \VUgras{перепелов} \textsuperscript{(42)}, \VUgras{дикую утку}
\textsuperscript{(43)} или \VUgras{фазана} \textsuperscript{(44)} --- как и
тропические плоды вроде \VUgras{бананов} \textsuperscript{(46)} и \VUgras{ананасов}.

%%K t15-09-1 p
9. --- Очень услужливый \VUgras{приказчик} \textsuperscript{(45)} готов подать
всё, что попросят: \VUgras{варенье} \textsuperscript{(47)}, смородиновое или
айвовое, \VUgras{смалец} \textsuperscript{(48)}, \VUgras{корнишоны}
\textsuperscript{(49)} или солёные \VUgras{сардины} \textsuperscript{(50)}.

%%K t15-09-2 p
Это очень удобно \VUgras{покупателям} \textsuperscript{(51)}, которые находят
рядом с самыми обыкновенными товарами редчайшие ранние овощи и самые
вкусные плоды: сочные \VUgras{груши} \textsuperscript{(52)}, \VUgras{сливы}
\textsuperscript{(53)}, \VUgras{виноград} \textsuperscript{(54)}, \VUgras{персики}
\textsuperscript{(55)}, \VUgras{миндаль} \textsuperscript{(56)} и \VUgras{орехи}
\textsuperscript{(57)}.

%%K t15-tit-5 sub
\VUsaut{4.0mm}
\VUpk{10.2pt}{40}{Покупатели.}
\VUsaut{3.0mm}

%%K t15-10-1 p
10. --- \VUgras{Рис} \textsuperscript{(58)} и \VUgras{чечевица} \textsuperscript{(59)}
стоят рядом с ящиком копчёных \VUgras{сельдей} \textsuperscript{(60)};
\VUgras{картофель} \textsuperscript{(61)} и \VUgras{фасоль} \textsuperscript{(63)}
соседствуют с \VUgras{яблоками} \textsuperscript{(62)}, \VUgras{инжиром}
\textsuperscript{(64)}, \VUgras{черносливом} \textsuperscript{(65)} и
\VUgras{орешками} \textsuperscript{(66)}. Свежие \VUgras{яйца} \textsuperscript{(67)} и
\VUgras{маслины} \textsuperscript{(68)} тоже есть в этой лавке, так что
\VUgras{корзины} \textsuperscript{(70)} и \VUgras{сетки} \textsuperscript{(73)}
наполняются очень быстро.

%%K t15-11-1 p
11. --- \VUgras{Кофе} \textsuperscript{(72)} этой отличной фирмы заслуженно
славится среди \VUgras{служанок} \textsuperscript{(69)} и кухарок, потому что
старый \VUgras{бакалейщик} \textsuperscript{(71)} жарит его сам с величайшим
тщанием.

%%K t15-tit-6 sub
\VUsaut{4.0mm}
\VUpk{10.2pt}{40}{Что тут можно увидеть.}
\VUsaut{3.0mm}

%%K t15-12-1 p
12. --- Не все приходят на рынок за припасами. В живописном движении
улочки, что ведёт к навесу, видишь самый разный народ. Рядом с добродушным
\VUgras{боенским рабочим} \textsuperscript{(75)}, который толкает перед собою
\VUgras{тачку} \textsuperscript{(74)}, вот два солдата в отпуску, очень гордые
своими яркими мундирами: \VUgras{гусар} \textsuperscript{(76)} в синем
\VUgras{доломане} и \VUgras{кивере с султаном} и \VUgras{капрал зуавов} в
широких шароварах и с \VUgras{феской} на голове.

%%K t15-13-1 p
13. --- \VUgras{Булочник} \textsuperscript{(80)}, стоящий на пороге
\VUgras{булочной} \textsuperscript{(78)}, только что вымесил тесто для своего
\VUgras{хлеба} \textsuperscript{(79)} и вынул из печи \VUgras{рогалики}
\textsuperscript{(81)}, \VUgras{булочки} \textsuperscript{(82)} и \VUgras{караваи}
\textsuperscript{(83)}, что лежат в витрине.

%%K t15-14-1 p
14. --- Над булочной живёт искусная \VUgras{швея} \textsuperscript{(84)}, у
которой работает несколько \VUgras{вышивальщиц} \textsuperscript{(85)}. Рядом с
нею живёт \VUgras{прачка и гладильщица} \textsuperscript{(86)}, которая крахмалит
\VUgras{бельё} \textsuperscript{(88)}, вымыв и высушив его, и гладит его
\VUgras{утюгом} \textsuperscript{(87)}.

%%K t15-tit-7 sub
\VUsaut{4.0mm}
\VUpk{10.2pt}{40}{Мясо, рыба и овощи.}
\VUsaut{3.0mm}

%%K t15-15-1 p
15. --- В \VUgras{мясной лавке} \textsuperscript{(89)} на нижнем этаже продают
всякое мясо --- \VUgras{баранину} \textsuperscript{(90)}, \VUgras{говядину}
\textsuperscript{(94)} или \VUgras{телятину} \textsuperscript{(92)} --- в виде
\VUgras{бараньих ног, вырезки, жаркого} и \VUgras{отбивных}.
\VUgras{Мясницкий подручный} \textsuperscript{(93)} разделывает всё это мясо и
откладывает \VUgras{мозговые кости} \textsuperscript{(96)}. У дверей мясной стоит
\VUgras{торговка устрицами} \textsuperscript{(97)}, которая продаёт
\VUgras{устрицы} \textsuperscript{(98)}. Она их открывает, если попросят.

%%K t15-16-1 p
16. --- Все эти торговцы платят патент или место́вое; вот почему
\VUgras{городовой} \textsuperscript{(100)} без пощады гоняет бедных
\VUgras{разносчиц овощей} \textsuperscript{(99)}, которые составляют им
конкуренцию, развозя на своих тележках \VUgras{морковь, редис, репу, лук-порей}
или \VUgras{пучки спаржи, свежий горошек, шпинат, щавель, кресс} и
\VUgras{петрушку}.

%%K t15-tit-8 sub
\VUsaut{4.0mm}
\VUpk{10.2pt}{40}{Берегись городового!}
\VUsaut{3.0mm}

%%K t15-17-1 p
17. --- Мальчик, торгующий \VUgras{чесноком} \textsuperscript{(101)} и
\VUgras{луком}, прекрасно знает, что и ему не позволено продавать свой
товар возле рынка. Он пускается бежать, рискуя опрокинуть корзину
\VUgras{крабов}, \VUgras{раков} и \VUgras{креветок}, что принадлежит
\VUgras{торговке} \textsuperscript{(102)} \VUgras{лангустами} и
\VUgras{омарами}.

%%K t15-18-1 p
18. --- Все хлопочут и поднимают большой шум; но это не мешает ясно
слышать зазыв бродячего \VUgras{стекольщика} \textsuperscript{(103)} и крик
\VUgras{угольщика} \textsuperscript{(104)}, который только что оставил на минуту
свою тележку, чтобы доставить \VUgras{мешок угля} \textsuperscript{(105)}.
\VUsaut{6.0mm}
\VUfilet{22mm}
